\section{Presentación del parcial}

\begin{objetivo}
Comprender el razonamiento computacional y diseñar algoritmos mediante
pseudocódigo, diagramas de flujo y pruebas de escritorio, utilizando
estructuras secuenciales, condicionales e iterativas.
\end{objetivo}

\subsection{Instrucciones generales}

\begin{enumerate}[leftmargin=*,label=\arabic*.]
  \item Lee cada actividad antes de responder.
  \item Conserva el procedimiento, no solamente el resultado final.
  \item Usa diagramas de flujo cuando la actividad lo indique.
  \item Anexa las evidencias solicitadas y revisa la rúbrica correspondiente.
\end{enumerate}

\subsection{Temario}

\begin{itemize}[leftmargin=*]
  \item Razonamiento computacional, algoritmos y granularidad.
  \item Variables, tipos de datos, operadores y nomenclatura.
  \item Diagramas de flujo, pseudocódigo y pruebas de escritorio.
  \item Estructuras condicionales e iterativas.
  \item Clasificación de errores y resolución de problemas integradores.
\end{itemize}
